% Figure from MATH 551 notes, section 13. Compile with the notes preamble.
\begin{tikzpicture}[x=1.0cm, y=1.0cm]
  \foreach \k/\y in {1/0,2/-1,3/-2,4/-3} {
    \draw[-stealth, thin] (-0.1,\y) -- (6.6,\y);
    \fill[figB!15] (\k,\y) rectangle (\k+1,\y+0.6);
    \draw[figB, thick] (-0.1,\y) -- (\k,\y) (\k+1,\y) -- (6.4,\y);
    \draw[figB, thick] (\k,\y+0.6) -- (\k+1,\y+0.6);
    \fill[figB] (\k,\y+0.6) circle (1.2pt) (\k+1,\y+0.6) circle (1.2pt);
    \filldraw[fill=white, draw=figB, thin] (\k,\y) circle (1.2pt) (\k+1,\y) circle (1.2pt);
    \node[font=\scriptsize, anchor=east] at (-0.15,\y+0.25) {$f_{\k}$};
  }
  \foreach \k in {1,...,5} \node[font=\scriptsize, black!70, below] at (\k,-3) {$\k$};
  \node[font=\scriptsize, black!60] at (3,-3.75) {$\vdots$};
  % a fixed point x: f_k(x) = 0 once k > x
  \draw[figR, densely dotted] (2.4,0.75) -- (2.4,-3.1);
  \node[font=\scriptsize, figR, anchor=south] at (2.4,0.75) {$x$};
  \foreach \y in {0,-2,-3} \fill[figR] (2.4,\y) circle (1.3pt);
  \fill[figR] (2.4,-0.4) circle (1.3pt);
  \draw[figR, -{Stealth[length=5pt]}, thick] (4.4,-3.55) -- (6.2,-3.55);
  \node[font=\scriptsize, figR, anchor=west] at (6.2,-3.55) {bump escapes};
\end{tikzpicture}
